% =============================================================================
% Chapter 2 addendum - High-level data pipeline overview
% =============================================================================
% BEST WAY TO USE THIS FILE: upload it to Overleaf as
% contents/chapter2_pipeline_overview.tex and add ONE line at the end of the
% background chapter, after its \section{Conclusion}:
%
%     % =============================================================================
% Chapter 2 addendum - High-level data pipeline overview
% =============================================================================
% BEST WAY TO USE THIS FILE: upload it to Overleaf as
% contents/chapter2_pipeline_overview.tex and add ONE line at the end of the
% background chapter, after its \section{Conclusion}:
%
%     % =============================================================================
% Chapter 2 addendum - High-level data pipeline overview
% =============================================================================
% BEST WAY TO USE THIS FILE: upload it to Overleaf as
% contents/chapter2_pipeline_overview.tex and add ONE line at the end of the
% background chapter, after its \section{Conclusion}:
%
%     % =============================================================================
% Chapter 2 addendum - High-level data pipeline overview
% =============================================================================
% BEST WAY TO USE THIS FILE: upload it to Overleaf as
% contents/chapter2_pipeline_overview.tex and add ONE line at the end of the
% background chapter, after its \section{Conclusion}:
%
%     \input{contents/chapter2_pipeline_overview}
%
% Do NOT copy-paste the contents into the chapter through a browser. A previous
% paste replaced the whole shape-trace table with a literal "<br>", which
% renders as the glyph pair "ibri" in T1 encoding. \input avoids that entirely.
%
% REQUIRED IN THE MASTER PREAMBLE (already present in this project's main.tex):
%   \usepackage{tikz}
%   \usetikzlibrary{arrows.meta, positioning, fit, backgrounds}
%
% The table uses \hline, matching every other table in the thesis, so no
% booktabs dependency is introduced. The figure is greyscale-safe: no meaning
% is carried by colour alone, since the switched components are also dashed.
%
% Labels introduced here: sec:pipeline-overview, fig:pipeline-overview,
% tab:pipeline-shapes, sec:pipeline-in-words. This file references no label
% outside itself, so it is safe to drop into either chapter2.tex or
% chapter2_shortened.tex.
% =============================================================================

\section{The Pipeline End to End}
\label{sec:pipeline-overview}

The preceding sections introduced each mechanism separately. This section puts
them in order, so that the rest of the thesis can refer to a single picture of
what happens to a clip between the corpus and a score.

The pipeline runs in \textbf{two stages}, and the division matters because the
two are executed a very different number of times.
\textbf{Stage~1 is preprocessing}: it converts each recorded clip into one
motion tensor, once, and saves it to disk. \textbf{Stage~2 is training and
evaluation}: it loads those tensors and fits a model, repeated for every
configuration and every held-out subject. Because Stage~1 is fixed before any
model is trained, every architectural comparison in this thesis is made on
identical inputs.

\autoref{fig:pipeline-overview} gives the whole sequence.
\autoref{tab:pipeline-shapes} traces what the data actually \emph{is} at each
step, which is often the quickest way to read a pipeline.

% -----------------------------------------------------------------------------
\begin{figure}[htbp]
	\centering
	\begin{tikzpicture}[
		font=\scriptsize,
		node distance=3.2mm,
		box/.style={
			draw, rounded corners=2pt, align=center,
			text width=34mm, inner sep=3pt,
			fill=black!4},
		switch/.style={box, dashed, fill=black!10},
		term/.style={box, fill=black!14, font=\scriptsize\bfseries},
		arr/.style={-{Stealth[length=1.6mm]}, semithick}
	]

	% ---------------- Stage 1: preprocessing ----------------
	\node[term] (a1) {CASME II corpus\\\normalfont 200\,fps, FACS-coded};
	\node[box, below=of a1] (a2)
		{Clip selection\\\normalfont 156 clips, 25 subjects};
	\node[box, below=of a2] (a3)
		{Frame preparation\\\normalfont greyscale, $224\times224$};
	\node[box, below=of a3] (a4)
		{Uniform sampling\\\normalfont 33 frames per clip};
	\node[switch, below=of a4] (a5)
		{Motion magnification\\\normalfont Eulerian, $\alpha=10$};
	\node[box, below=of a5] (a6)
		{Optical flow\\\normalfont Farneb\"ack, 32 pairs};
	\node[box, below=of a6] (a7)
		{Optical strain\\\normalfont spatial gradients};
	\node[box, below=of a7] (a8)
		{Stack and normalise\\\normalfont $[u,v,\varepsilon]$ to $[0,1]$};
	\node[term, below=of a8] (a9)
		{Saved motion tensor\\\normalfont $(3,32,224,224)$};

	\foreach \i/\j in {a1/a2, a2/a3, a3/a4, a4/a5, a5/a6, a6/a7, a7/a8, a8/a9}
		{\draw[arr] (\i) -- (\j);}

	% ---------------- Stage 2: training and evaluation ----------------
	\node[box, right=20mm of a2] (b1)
		{Load and standardise\\\normalfont flip augmentation};
	\node[switch, below=of b1] (b2)
		{Convolutional stem\\\normalfont three streams};
	\node[switch, below=of b2] (b3)
		{SimAM\\\normalfont parameter-free gate};
	\node[box, below=of b3] (b4)
		{Pool to a sequence\\\normalfont 32 steps, 96 features};
	\node[switch, below=of b4] (b5)
		{Temporal transformer\\\normalfont 4 layers, 8 heads};
	\node[box, below=of b5] (b6)
		{Classifier head\\\normalfont three logits};
	\node[box, below=of b6] (b7)
		{Leave-one-subject-out\\\normalfont 25 folds, counts pooled};
	\node[term, below=of b7] (b8)
		{Pooled macro F1\\\normalfont the primary metric};

	\foreach \i/\j in {b1/b2, b2/b3, b3/b4, b4/b5, b5/b6, b6/b7, b7/b8}
		{\draw[arr] (\i) -- (\j);}

	% ---------------- the hand-off ----------------
	\draw[arr] (a9.east) -- ++(8mm,0) |- (b1.west);

	% ---------------- stage brackets ----------------
	\begin{scope}[on background layer]
		\node[draw=black!45, dashed, rounded corners=3pt, inner sep=4.5mm,
		      fit=(a1)(a9)] (s1) {};
		\node[draw=black!45, dashed, rounded corners=3pt, inner sep=4.5mm,
		      fit=(b1)(b8)] (s2) {};
	\end{scope}
	\node[above=0.5mm of s1, font=\scriptsize\bfseries] {Stage 1 --- preprocessing};
	\node[above=0.5mm of s2, font=\scriptsize\bfseries] {Stage 2 --- training and evaluation};

	\end{tikzpicture}
	\caption[The data pipeline, end to end]{The pipeline from corpus to score.
	Stage~1 runs once per clip; Stage~2 runs once per configuration, per fold.
	Solid boxes are always executed; \textbf{dashed boxes are the four
	components the ablation switches on and off}, each of which has a defined
	replacement when disabled rather than an empty slot. Because Stage~1 is
	fixed before any model is trained, every architectural comparison uses
	identical inputs.}
	\label{fig:pipeline-overview}
\end{figure}
% -----------------------------------------------------------------------------

\subsection{The sequence in words}
\label{sec:pipeline-in-words}

\textbf{Stage 1 turns video into motion.} Clips are drawn from CASME~II, kept
if they are coded as micro-expressions and carry a label in the three-class
mapping, and reduced to \textbf{156 clips across 25 subjects}. Each is loaded
in greyscale, resized to $224\times224$, and sampled at 33 evenly spaced
frames between the annotated onset and offset. Magnification is applied at
this point if enabled. Optical flow is then computed for the 32 adjacent frame
pairs, giving horizontal and vertical displacement fields; optical strain is
derived from their spatial gradients. The three are stacked and normalised
into one saved tensor.

The important consequence is that \textbf{the network never sees a pixel
intensity sequence}. It sees displacement and deformation, computed
analytically before training begins.

\textbf{Stage 2 turns motion into a score.} A saved tensor is loaded,
optionally flipped as augmentation, and standardised per channel. It passes
through a convolutional stem, an attention gate, a pooling step that collapses
space and leaves time, a temporal encoder, and a three-way classifier head.
Training uses balanced sampling and a focal objective with label smoothing.
Scoring is leave-one-subject-out over all 25 subjects, with confusion counts
accumulated across folds before any per-class score is computed.

\subsection{What the data is at each step}
\label{sec:pipeline-shapes-discussion}

\begin{table}[htbp]
	\centering
	\small
	\begin{tabular}{@{}lll@{}}
		\hline
		\textbf{After this step} & \textbf{The data is} & \textbf{Shape} \\
		\hline
		Corpus selection      & a set of clips                & 156 clips, 25 subjects \\
		Frame preparation     & greyscale frames              & $N \times 224 \times 224$ \\
		Uniform sampling      & a fixed-length sequence       & $33 \times 224 \times 224$ \\
		Magnification         & the same, amplified           & $33 \times 224 \times 224$ \\
		Flow and strain       & three motion fields per pair  & $3 \times 32 \times 224 \times 224$ \\
		Pooling to a sequence & a sequence of feature vectors & $32 \times 96$ \\
		Temporal encoding     & one vector per clip           & $96$ \\
		Classifier head       & class scores                  & $3$ \\
		Fold aggregation      & pooled confusion counts       & $3 \times 3$ \\
		\hline
	\end{tabular}
	\caption[What the data is at each pipeline step]{The representation at each
	stage. The temporal axis falls from 33 frames to 32 pairs because flow is
	computed between \emph{adjacent} frames, and the spatial axes vanish at the
	pooling step rather than at the classifier.}
	\label{tab:pipeline-shapes}
\end{table}

\subsection{Where the experiment intervenes}
\label{sec:pipeline-intervenes}

Four of the boxes in \autoref{fig:pipeline-overview} are dashed because they
are the subject of this thesis rather than fixed infrastructure: motion
magnification, the convolutional stem, SimAM and the temporal transformer.
Each can be switched off independently, and each has a defined replacement
when it is --- not an empty slot. The component-controls table earlier in this
chapter states each replacement. Everything else in the diagram is held
identical across every configuration and every fold.

That is what makes the comparison a controlled one, and it is the design whose
results Chapter~5 reports.

%
% Do NOT copy-paste the contents into the chapter through a browser. A previous
% paste replaced the whole shape-trace table with a literal "<br>", which
% renders as the glyph pair "ibri" in T1 encoding. \input avoids that entirely.
%
% REQUIRED IN THE MASTER PREAMBLE (already present in this project's main.tex):
%   \usepackage{tikz}
%   \usetikzlibrary{arrows.meta, positioning, fit, backgrounds}
%
% The table uses \hline, matching every other table in the thesis, so no
% booktabs dependency is introduced. The figure is greyscale-safe: no meaning
% is carried by colour alone, since the switched components are also dashed.
%
% Labels introduced here: sec:pipeline-overview, fig:pipeline-overview,
% tab:pipeline-shapes, sec:pipeline-in-words. This file references no label
% outside itself, so it is safe to drop into either chapter2.tex or
% chapter2_shortened.tex.
% =============================================================================

\section{The Pipeline End to End}
\label{sec:pipeline-overview}

The preceding sections introduced each mechanism separately. This section puts
them in order, so that the rest of the thesis can refer to a single picture of
what happens to a clip between the corpus and a score.

The pipeline runs in \textbf{two stages}, and the division matters because the
two are executed a very different number of times.
\textbf{Stage~1 is preprocessing}: it converts each recorded clip into one
motion tensor, once, and saves it to disk. \textbf{Stage~2 is training and
evaluation}: it loads those tensors and fits a model, repeated for every
configuration and every held-out subject. Because Stage~1 is fixed before any
model is trained, every architectural comparison in this thesis is made on
identical inputs.

\autoref{fig:pipeline-overview} gives the whole sequence.
\autoref{tab:pipeline-shapes} traces what the data actually \emph{is} at each
step, which is often the quickest way to read a pipeline.

% -----------------------------------------------------------------------------
\begin{figure}[htbp]
	\centering
	\begin{tikzpicture}[
		font=\scriptsize,
		node distance=3.2mm,
		box/.style={
			draw, rounded corners=2pt, align=center,
			text width=34mm, inner sep=3pt,
			fill=black!4},
		switch/.style={box, dashed, fill=black!10},
		term/.style={box, fill=black!14, font=\scriptsize\bfseries},
		arr/.style={-{Stealth[length=1.6mm]}, semithick}
	]

	% ---------------- Stage 1: preprocessing ----------------
	\node[term] (a1) {CASME II corpus\\\normalfont 200\,fps, FACS-coded};
	\node[box, below=of a1] (a2)
		{Clip selection\\\normalfont 156 clips, 25 subjects};
	\node[box, below=of a2] (a3)
		{Frame preparation\\\normalfont greyscale, $224\times224$};
	\node[box, below=of a3] (a4)
		{Uniform sampling\\\normalfont 33 frames per clip};
	\node[switch, below=of a4] (a5)
		{Motion magnification\\\normalfont Eulerian, $\alpha=10$};
	\node[box, below=of a5] (a6)
		{Optical flow\\\normalfont Farneb\"ack, 32 pairs};
	\node[box, below=of a6] (a7)
		{Optical strain\\\normalfont spatial gradients};
	\node[box, below=of a7] (a8)
		{Stack and normalise\\\normalfont $[u,v,\varepsilon]$ to $[0,1]$};
	\node[term, below=of a8] (a9)
		{Saved motion tensor\\\normalfont $(3,32,224,224)$};

	\foreach \i/\j in {a1/a2, a2/a3, a3/a4, a4/a5, a5/a6, a6/a7, a7/a8, a8/a9}
		{\draw[arr] (\i) -- (\j);}

	% ---------------- Stage 2: training and evaluation ----------------
	\node[box, right=20mm of a2] (b1)
		{Load and standardise\\\normalfont flip augmentation};
	\node[switch, below=of b1] (b2)
		{Convolutional stem\\\normalfont three streams};
	\node[switch, below=of b2] (b3)
		{SimAM\\\normalfont parameter-free gate};
	\node[box, below=of b3] (b4)
		{Pool to a sequence\\\normalfont 32 steps, 96 features};
	\node[switch, below=of b4] (b5)
		{Temporal transformer\\\normalfont 4 layers, 8 heads};
	\node[box, below=of b5] (b6)
		{Classifier head\\\normalfont three logits};
	\node[box, below=of b6] (b7)
		{Leave-one-subject-out\\\normalfont 25 folds, counts pooled};
	\node[term, below=of b7] (b8)
		{Pooled macro F1\\\normalfont the primary metric};

	\foreach \i/\j in {b1/b2, b2/b3, b3/b4, b4/b5, b5/b6, b6/b7, b7/b8}
		{\draw[arr] (\i) -- (\j);}

	% ---------------- the hand-off ----------------
	\draw[arr] (a9.east) -- ++(8mm,0) |- (b1.west);

	% ---------------- stage brackets ----------------
	\begin{scope}[on background layer]
		\node[draw=black!45, dashed, rounded corners=3pt, inner sep=4.5mm,
		      fit=(a1)(a9)] (s1) {};
		\node[draw=black!45, dashed, rounded corners=3pt, inner sep=4.5mm,
		      fit=(b1)(b8)] (s2) {};
	\end{scope}
	\node[above=0.5mm of s1, font=\scriptsize\bfseries] {Stage 1 --- preprocessing};
	\node[above=0.5mm of s2, font=\scriptsize\bfseries] {Stage 2 --- training and evaluation};

	\end{tikzpicture}
	\caption[The data pipeline, end to end]{The pipeline from corpus to score.
	Stage~1 runs once per clip; Stage~2 runs once per configuration, per fold.
	Solid boxes are always executed; \textbf{dashed boxes are the four
	components the ablation switches on and off}, each of which has a defined
	replacement when disabled rather than an empty slot. Because Stage~1 is
	fixed before any model is trained, every architectural comparison uses
	identical inputs.}
	\label{fig:pipeline-overview}
\end{figure}
% -----------------------------------------------------------------------------

\subsection{The sequence in words}
\label{sec:pipeline-in-words}

\textbf{Stage 1 turns video into motion.} Clips are drawn from CASME~II, kept
if they are coded as micro-expressions and carry a label in the three-class
mapping, and reduced to \textbf{156 clips across 25 subjects}. Each is loaded
in greyscale, resized to $224\times224$, and sampled at 33 evenly spaced
frames between the annotated onset and offset. Magnification is applied at
this point if enabled. Optical flow is then computed for the 32 adjacent frame
pairs, giving horizontal and vertical displacement fields; optical strain is
derived from their spatial gradients. The three are stacked and normalised
into one saved tensor.

The important consequence is that \textbf{the network never sees a pixel
intensity sequence}. It sees displacement and deformation, computed
analytically before training begins.

\textbf{Stage 2 turns motion into a score.} A saved tensor is loaded,
optionally flipped as augmentation, and standardised per channel. It passes
through a convolutional stem, an attention gate, a pooling step that collapses
space and leaves time, a temporal encoder, and a three-way classifier head.
Training uses balanced sampling and a focal objective with label smoothing.
Scoring is leave-one-subject-out over all 25 subjects, with confusion counts
accumulated across folds before any per-class score is computed.

\subsection{What the data is at each step}
\label{sec:pipeline-shapes-discussion}

\begin{table}[htbp]
	\centering
	\small
	\begin{tabular}{@{}lll@{}}
		\hline
		\textbf{After this step} & \textbf{The data is} & \textbf{Shape} \\
		\hline
		Corpus selection      & a set of clips                & 156 clips, 25 subjects \\
		Frame preparation     & greyscale frames              & $N \times 224 \times 224$ \\
		Uniform sampling      & a fixed-length sequence       & $33 \times 224 \times 224$ \\
		Magnification         & the same, amplified           & $33 \times 224 \times 224$ \\
		Flow and strain       & three motion fields per pair  & $3 \times 32 \times 224 \times 224$ \\
		Pooling to a sequence & a sequence of feature vectors & $32 \times 96$ \\
		Temporal encoding     & one vector per clip           & $96$ \\
		Classifier head       & class scores                  & $3$ \\
		Fold aggregation      & pooled confusion counts       & $3 \times 3$ \\
		\hline
	\end{tabular}
	\caption[What the data is at each pipeline step]{The representation at each
	stage. The temporal axis falls from 33 frames to 32 pairs because flow is
	computed between \emph{adjacent} frames, and the spatial axes vanish at the
	pooling step rather than at the classifier.}
	\label{tab:pipeline-shapes}
\end{table}

\subsection{Where the experiment intervenes}
\label{sec:pipeline-intervenes}

Four of the boxes in \autoref{fig:pipeline-overview} are dashed because they
are the subject of this thesis rather than fixed infrastructure: motion
magnification, the convolutional stem, SimAM and the temporal transformer.
Each can be switched off independently, and each has a defined replacement
when it is --- not an empty slot. The component-controls table earlier in this
chapter states each replacement. Everything else in the diagram is held
identical across every configuration and every fold.

That is what makes the comparison a controlled one, and it is the design whose
results Chapter~5 reports.

%
% Do NOT copy-paste the contents into the chapter through a browser. A previous
% paste replaced the whole shape-trace table with a literal "<br>", which
% renders as the glyph pair "ibri" in T1 encoding. \input avoids that entirely.
%
% REQUIRED IN THE MASTER PREAMBLE (already present in this project's main.tex):
%   \usepackage{tikz}
%   \usetikzlibrary{arrows.meta, positioning, fit, backgrounds}
%
% The table uses \hline, matching every other table in the thesis, so no
% booktabs dependency is introduced. The figure is greyscale-safe: no meaning
% is carried by colour alone, since the switched components are also dashed.
%
% Labels introduced here: sec:pipeline-overview, fig:pipeline-overview,
% tab:pipeline-shapes, sec:pipeline-in-words. This file references no label
% outside itself, so it is safe to drop into either chapter2.tex or
% chapter2_shortened.tex.
% =============================================================================

\section{The Pipeline End to End}
\label{sec:pipeline-overview}

The preceding sections introduced each mechanism separately. This section puts
them in order, so that the rest of the thesis can refer to a single picture of
what happens to a clip between the corpus and a score.

The pipeline runs in \textbf{two stages}, and the division matters because the
two are executed a very different number of times.
\textbf{Stage~1 is preprocessing}: it converts each recorded clip into one
motion tensor, once, and saves it to disk. \textbf{Stage~2 is training and
evaluation}: it loads those tensors and fits a model, repeated for every
configuration and every held-out subject. Because Stage~1 is fixed before any
model is trained, every architectural comparison in this thesis is made on
identical inputs.

\autoref{fig:pipeline-overview} gives the whole sequence.
\autoref{tab:pipeline-shapes} traces what the data actually \emph{is} at each
step, which is often the quickest way to read a pipeline.

% -----------------------------------------------------------------------------
\begin{figure}[htbp]
	\centering
	\begin{tikzpicture}[
		font=\scriptsize,
		node distance=3.2mm,
		box/.style={
			draw, rounded corners=2pt, align=center,
			text width=34mm, inner sep=3pt,
			fill=black!4},
		switch/.style={box, dashed, fill=black!10},
		term/.style={box, fill=black!14, font=\scriptsize\bfseries},
		arr/.style={-{Stealth[length=1.6mm]}, semithick}
	]

	% ---------------- Stage 1: preprocessing ----------------
	\node[term] (a1) {CASME II corpus\\\normalfont 200\,fps, FACS-coded};
	\node[box, below=of a1] (a2)
		{Clip selection\\\normalfont 156 clips, 25 subjects};
	\node[box, below=of a2] (a3)
		{Frame preparation\\\normalfont greyscale, $224\times224$};
	\node[box, below=of a3] (a4)
		{Uniform sampling\\\normalfont 33 frames per clip};
	\node[switch, below=of a4] (a5)
		{Motion magnification\\\normalfont Eulerian, $\alpha=10$};
	\node[box, below=of a5] (a6)
		{Optical flow\\\normalfont Farneb\"ack, 32 pairs};
	\node[box, below=of a6] (a7)
		{Optical strain\\\normalfont spatial gradients};
	\node[box, below=of a7] (a8)
		{Stack and normalise\\\normalfont $[u,v,\varepsilon]$ to $[0,1]$};
	\node[term, below=of a8] (a9)
		{Saved motion tensor\\\normalfont $(3,32,224,224)$};

	\foreach \i/\j in {a1/a2, a2/a3, a3/a4, a4/a5, a5/a6, a6/a7, a7/a8, a8/a9}
		{\draw[arr] (\i) -- (\j);}

	% ---------------- Stage 2: training and evaluation ----------------
	\node[box, right=20mm of a2] (b1)
		{Load and standardise\\\normalfont flip augmentation};
	\node[switch, below=of b1] (b2)
		{Convolutional stem\\\normalfont three streams};
	\node[switch, below=of b2] (b3)
		{SimAM\\\normalfont parameter-free gate};
	\node[box, below=of b3] (b4)
		{Pool to a sequence\\\normalfont 32 steps, 96 features};
	\node[switch, below=of b4] (b5)
		{Temporal transformer\\\normalfont 4 layers, 8 heads};
	\node[box, below=of b5] (b6)
		{Classifier head\\\normalfont three logits};
	\node[box, below=of b6] (b7)
		{Leave-one-subject-out\\\normalfont 25 folds, counts pooled};
	\node[term, below=of b7] (b8)
		{Pooled macro F1\\\normalfont the primary metric};

	\foreach \i/\j in {b1/b2, b2/b3, b3/b4, b4/b5, b5/b6, b6/b7, b7/b8}
		{\draw[arr] (\i) -- (\j);}

	% ---------------- the hand-off ----------------
	\draw[arr] (a9.east) -- ++(8mm,0) |- (b1.west);

	% ---------------- stage brackets ----------------
	\begin{scope}[on background layer]
		\node[draw=black!45, dashed, rounded corners=3pt, inner sep=4.5mm,
		      fit=(a1)(a9)] (s1) {};
		\node[draw=black!45, dashed, rounded corners=3pt, inner sep=4.5mm,
		      fit=(b1)(b8)] (s2) {};
	\end{scope}
	\node[above=0.5mm of s1, font=\scriptsize\bfseries] {Stage 1 --- preprocessing};
	\node[above=0.5mm of s2, font=\scriptsize\bfseries] {Stage 2 --- training and evaluation};

	\end{tikzpicture}
	\caption[The data pipeline, end to end]{The pipeline from corpus to score.
	Stage~1 runs once per clip; Stage~2 runs once per configuration, per fold.
	Solid boxes are always executed; \textbf{dashed boxes are the four
	components the ablation switches on and off}, each of which has a defined
	replacement when disabled rather than an empty slot. Because Stage~1 is
	fixed before any model is trained, every architectural comparison uses
	identical inputs.}
	\label{fig:pipeline-overview}
\end{figure}
% -----------------------------------------------------------------------------

\subsection{The sequence in words}
\label{sec:pipeline-in-words}

\textbf{Stage 1 turns video into motion.} Clips are drawn from CASME~II, kept
if they are coded as micro-expressions and carry a label in the three-class
mapping, and reduced to \textbf{156 clips across 25 subjects}. Each is loaded
in greyscale, resized to $224\times224$, and sampled at 33 evenly spaced
frames between the annotated onset and offset. Magnification is applied at
this point if enabled. Optical flow is then computed for the 32 adjacent frame
pairs, giving horizontal and vertical displacement fields; optical strain is
derived from their spatial gradients. The three are stacked and normalised
into one saved tensor.

The important consequence is that \textbf{the network never sees a pixel
intensity sequence}. It sees displacement and deformation, computed
analytically before training begins.

\textbf{Stage 2 turns motion into a score.} A saved tensor is loaded,
optionally flipped as augmentation, and standardised per channel. It passes
through a convolutional stem, an attention gate, a pooling step that collapses
space and leaves time, a temporal encoder, and a three-way classifier head.
Training uses balanced sampling and a focal objective with label smoothing.
Scoring is leave-one-subject-out over all 25 subjects, with confusion counts
accumulated across folds before any per-class score is computed.

\subsection{What the data is at each step}
\label{sec:pipeline-shapes-discussion}

\begin{table}[htbp]
	\centering
	\small
	\begin{tabular}{@{}lll@{}}
		\hline
		\textbf{After this step} & \textbf{The data is} & \textbf{Shape} \\
		\hline
		Corpus selection      & a set of clips                & 156 clips, 25 subjects \\
		Frame preparation     & greyscale frames              & $N \times 224 \times 224$ \\
		Uniform sampling      & a fixed-length sequence       & $33 \times 224 \times 224$ \\
		Magnification         & the same, amplified           & $33 \times 224 \times 224$ \\
		Flow and strain       & three motion fields per pair  & $3 \times 32 \times 224 \times 224$ \\
		Pooling to a sequence & a sequence of feature vectors & $32 \times 96$ \\
		Temporal encoding     & one vector per clip           & $96$ \\
		Classifier head       & class scores                  & $3$ \\
		Fold aggregation      & pooled confusion counts       & $3 \times 3$ \\
		\hline
	\end{tabular}
	\caption[What the data is at each pipeline step]{The representation at each
	stage. The temporal axis falls from 33 frames to 32 pairs because flow is
	computed between \emph{adjacent} frames, and the spatial axes vanish at the
	pooling step rather than at the classifier.}
	\label{tab:pipeline-shapes}
\end{table}

\subsection{Where the experiment intervenes}
\label{sec:pipeline-intervenes}

Four of the boxes in \autoref{fig:pipeline-overview} are dashed because they
are the subject of this thesis rather than fixed infrastructure: motion
magnification, the convolutional stem, SimAM and the temporal transformer.
Each can be switched off independently, and each has a defined replacement
when it is --- not an empty slot. The component-controls table earlier in this
chapter states each replacement. Everything else in the diagram is held
identical across every configuration and every fold.

That is what makes the comparison a controlled one, and it is the design whose
results Chapter~5 reports.

%
% Do NOT copy-paste the contents into the chapter through a browser. A previous
% paste replaced the whole shape-trace table with a literal "<br>", which
% renders as the glyph pair "ibri" in T1 encoding. \input avoids that entirely.
%
% REQUIRED IN THE MASTER PREAMBLE (already present in this project's main.tex):
%   \usepackage{tikz}
%   \usetikzlibrary{arrows.meta, positioning, fit, backgrounds}
%
% The table uses \hline, matching every other table in the thesis, so no
% booktabs dependency is introduced. The figure is greyscale-safe: no meaning
% is carried by colour alone, since the switched components are also dashed.
%
% Labels introduced here: sec:pipeline-overview, fig:pipeline-overview,
% tab:pipeline-shapes, sec:pipeline-in-words. This file references no label
% outside itself, so it is safe to drop into either chapter2.tex or
% chapter2_shortened.tex.
% =============================================================================

\section{The Pipeline End to End}
\label{sec:pipeline-overview}

The preceding sections introduced each mechanism separately. This section puts
them in order, so that the rest of the thesis can refer to a single picture of
what happens to a clip between the corpus and a score.

The pipeline runs in \textbf{two stages}, and the division matters because the
two are executed a very different number of times.
\textbf{Stage~1 is preprocessing}: it converts each recorded clip into one
motion tensor, once, and saves it to disk. \textbf{Stage~2 is training and
evaluation}: it loads those tensors and fits a model, repeated for every
configuration and every held-out subject. Because Stage~1 is fixed before any
model is trained, every architectural comparison in this thesis is made on
identical inputs.

\autoref{fig:pipeline-overview} gives the whole sequence.
\autoref{tab:pipeline-shapes} traces what the data actually \emph{is} at each
step, which is often the quickest way to read a pipeline.

% -----------------------------------------------------------------------------
\begin{figure}[htbp]
	\centering
	\begin{tikzpicture}[
		font=\scriptsize,
		node distance=3.2mm,
		box/.style={
			draw, rounded corners=2pt, align=center,
			text width=34mm, inner sep=3pt,
			fill=black!4},
		switch/.style={box, dashed, fill=black!10},
		term/.style={box, fill=black!14, font=\scriptsize\bfseries},
		arr/.style={-{Stealth[length=1.6mm]}, semithick}
	]

	% ---------------- Stage 1: preprocessing ----------------
	\node[term] (a1) {CASME II corpus\\\normalfont 200\,fps, FACS-coded};
	\node[box, below=of a1] (a2)
		{Clip selection\\\normalfont 156 clips, 25 subjects};
	\node[box, below=of a2] (a3)
		{Frame preparation\\\normalfont greyscale, $224\times224$};
	\node[box, below=of a3] (a4)
		{Uniform sampling\\\normalfont 33 frames per clip};
	\node[switch, below=of a4] (a5)
		{Motion magnification\\\normalfont Eulerian, $\alpha=10$};
	\node[box, below=of a5] (a6)
		{Optical flow\\\normalfont Farneb\"ack, 32 pairs};
	\node[box, below=of a6] (a7)
		{Optical strain\\\normalfont spatial gradients};
	\node[box, below=of a7] (a8)
		{Stack and normalise\\\normalfont $[u,v,\varepsilon]$ to $[0,1]$};
	\node[term, below=of a8] (a9)
		{Saved motion tensor\\\normalfont $(3,32,224,224)$};

	\foreach \i/\j in {a1/a2, a2/a3, a3/a4, a4/a5, a5/a6, a6/a7, a7/a8, a8/a9}
		{\draw[arr] (\i) -- (\j);}

	% ---------------- Stage 2: training and evaluation ----------------
	\node[box, right=20mm of a2] (b1)
		{Load and standardise\\\normalfont flip augmentation};
	\node[switch, below=of b1] (b2)
		{Convolutional stem\\\normalfont three streams};
	\node[switch, below=of b2] (b3)
		{SimAM\\\normalfont parameter-free gate};
	\node[box, below=of b3] (b4)
		{Pool to a sequence\\\normalfont 32 steps, 96 features};
	\node[switch, below=of b4] (b5)
		{Temporal transformer\\\normalfont 4 layers, 8 heads};
	\node[box, below=of b5] (b6)
		{Classifier head\\\normalfont three logits};
	\node[box, below=of b6] (b7)
		{Leave-one-subject-out\\\normalfont 25 folds, counts pooled};
	\node[term, below=of b7] (b8)
		{Pooled macro F1\\\normalfont the primary metric};

	\foreach \i/\j in {b1/b2, b2/b3, b3/b4, b4/b5, b5/b6, b6/b7, b7/b8}
		{\draw[arr] (\i) -- (\j);}

	% ---------------- the hand-off ----------------
	\draw[arr] (a9.east) -- ++(8mm,0) |- (b1.west);

	% ---------------- stage brackets ----------------
	\begin{scope}[on background layer]
		\node[draw=black!45, dashed, rounded corners=3pt, inner sep=4.5mm,
		      fit=(a1)(a9)] (s1) {};
		\node[draw=black!45, dashed, rounded corners=3pt, inner sep=4.5mm,
		      fit=(b1)(b8)] (s2) {};
	\end{scope}
	\node[above=0.5mm of s1, font=\scriptsize\bfseries] {Stage 1 --- preprocessing};
	\node[above=0.5mm of s2, font=\scriptsize\bfseries] {Stage 2 --- training and evaluation};

	\end{tikzpicture}
	\caption[The data pipeline, end to end]{The pipeline from corpus to score.
	Stage~1 runs once per clip; Stage~2 runs once per configuration, per fold.
	Solid boxes are always executed; \textbf{dashed boxes are the four
	components the ablation switches on and off}, each of which has a defined
	replacement when disabled rather than an empty slot. Because Stage~1 is
	fixed before any model is trained, every architectural comparison uses
	identical inputs.}
	\label{fig:pipeline-overview}
\end{figure}
% -----------------------------------------------------------------------------

\subsection{The sequence in words}
\label{sec:pipeline-in-words}

\textbf{Stage 1 turns video into motion.} Clips are drawn from CASME~II, kept
if they are coded as micro-expressions and carry a label in the three-class
mapping, and reduced to \textbf{156 clips across 25 subjects}. Each is loaded
in greyscale, resized to $224\times224$, and sampled at 33 evenly spaced
frames between the annotated onset and offset. Magnification is applied at
this point if enabled. Optical flow is then computed for the 32 adjacent frame
pairs, giving horizontal and vertical displacement fields; optical strain is
derived from their spatial gradients. The three are stacked and normalised
into one saved tensor.

The important consequence is that \textbf{the network never sees a pixel
intensity sequence}. It sees displacement and deformation, computed
analytically before training begins.

\textbf{Stage 2 turns motion into a score.} A saved tensor is loaded,
optionally flipped as augmentation, and standardised per channel. It passes
through a convolutional stem, an attention gate, a pooling step that collapses
space and leaves time, a temporal encoder, and a three-way classifier head.
Training uses balanced sampling and a focal objective with label smoothing.
Scoring is leave-one-subject-out over all 25 subjects, with confusion counts
accumulated across folds before any per-class score is computed.

\subsection{What the data is at each step}
\label{sec:pipeline-shapes-discussion}

\begin{table}[htbp]
	\centering
	\small
	\begin{tabular}{@{}lll@{}}
		\hline
		\textbf{After this step} & \textbf{The data is} & \textbf{Shape} \\
		\hline
		Corpus selection      & a set of clips                & 156 clips, 25 subjects \\
		Frame preparation     & greyscale frames              & $N \times 224 \times 224$ \\
		Uniform sampling      & a fixed-length sequence       & $33 \times 224 \times 224$ \\
		Magnification         & the same, amplified           & $33 \times 224 \times 224$ \\
		Flow and strain       & three motion fields per pair  & $3 \times 32 \times 224 \times 224$ \\
		Pooling to a sequence & a sequence of feature vectors & $32 \times 96$ \\
		Temporal encoding     & one vector per clip           & $96$ \\
		Classifier head       & class scores                  & $3$ \\
		Fold aggregation      & pooled confusion counts       & $3 \times 3$ \\
		\hline
	\end{tabular}
	\caption[What the data is at each pipeline step]{The representation at each
	stage. The temporal axis falls from 33 frames to 32 pairs because flow is
	computed between \emph{adjacent} frames, and the spatial axes vanish at the
	pooling step rather than at the classifier.}
	\label{tab:pipeline-shapes}
\end{table}

\subsection{Where the experiment intervenes}
\label{sec:pipeline-intervenes}

Four of the boxes in \autoref{fig:pipeline-overview} are dashed because they
are the subject of this thesis rather than fixed infrastructure: motion
magnification, the convolutional stem, SimAM and the temporal transformer.
Each can be switched off independently, and each has a defined replacement
when it is --- not an empty slot. The component-controls table earlier in this
chapter states each replacement. Everything else in the diagram is held
identical across every configuration and every fold.

That is what makes the comparison a controlled one, and it is the design whose
results Chapter~5 reports.
